% =====================================================================
\section{Акселерометр и гироскоп: датчик движения}
% =====================================================================
\lessonintro
  {читать датчик напрямую через его регистры, измерять наклон и убирать шум и дрейф фильтром}
  {урок 7 (шина I\textsuperscript{2}C), урок 22 (регистры, старший и младший байты), урок 6 (шум и усреднение)}
  {плата измеряет угол наклона и рисует в Serial Plotter «сырой» и отфильтрованный графики}
  {модуль MPU6050 (плата GY-521)}

\subsection{Сначала: два датчика, которые ошибаются по-разному}

В модуле MPU6050 два датчика. Акселерометр чувствует ускорение, в том числе силу тяжести, поэтому по
нему видно, куда направлен «низ». Гироскоп чувствует скорость поворота. Каждый из них по отдельности
плох: акселерометр дрожит от любой вибрации, а угол, вычисленный по гироскопу, постепенно «уплывает».
Но ошибки у них разные, и если правильно смешать показания, получится точный и спокойный угол. Тот же
приём стоит в каждом телефоне, квадрокоптере и самобалансирующемся роботе.

\subsection{Что измеряет датчик}

\begin{figure}[H]
\centering
\begin{tikzpicture}[font=\sffamily\scriptsize]
  % плата горизонтально
  \begin{scope}
    \fill[blkper,draw=black] (-1.4,-0.25) rectangle (1.4,0.25);
    \node at (0,0) {MPU6050};
    \draw[arr,esp] (0,0.25) -- (0,1.5) node[above]{Z};
    \draw[arr,accent] (1.4,0) -- (2.5,0) node[right]{Y};
    \draw[arr,okgreen] (-0.5,-0.25) -- (-1.2,-0.95) node[below left]{X};
    \draw[arr,very thick] (0,-0.25) -- (0,-1.5) node[below]{сила тяжести $g$};
    \node[align=center] at (0,-2.6) {плата лежит ровно:\\ $a_z = 1g$, $a_y = 0$, угол 0°};
  \end{scope}
  % плата с наклоном
  \begin{scope}[shift={(7.5,0)}]
    \begin{scope}[rotate=30]
      \fill[blkper,draw=black] (-1.4,-0.25) rectangle (1.4,0.25);
      \draw[arr,esp] (0,0.25) -- (0,1.5) node[above]{Z};
      \draw[arr,accent] (1.4,0) -- (2.5,0) node[right]{Y};
    \end{scope}
    \draw[arr,very thick] (0,0) -- (0,-1.5) node[below]{$g$};
    \draw[dashed] (0,0) -- (2.4,0);
    \draw[->] (1.6,0) arc (0:30:1.6) node[midway,right]{$\varphi$};
    \node[align=center] at (0.6,-2.6) {плата наклонена вокруг оси X:\\ $\varphi = \operatorname{arctg}(a_y / a_z)$};
  \end{scope}
\end{tikzpicture}
\caption{Оси датчика и измерение наклона. Сила тяжести раскладывается по осям Y и Z, и по их отношению находят угол}
\end{figure}

\begin{center}\small
\begin{tabularx}{\textwidth}{@{}l X X@{}}
\toprule
 & Акселерометр & Гироскоп \\
\midrule
Что измеряет & ускорение по трём осям (в долях $g$) & скорость поворота вокруг трёх осей (°/с) \\
Как получить угол & сразу, по направлению силы тяжести & накапливая: угол += скорость × время \\
Сильная сторона & не накапливает ошибку & не чувствует вибрацию, быстрый \\
Слабая сторона & шумит и «врёт» при рывках & угол постепенно уплывает (дрейф) \\
\bottomrule
\end{tabularx}
\end{center}

\subsection{Разговор с датчиком через регистры}

В уроке 7 всю работу с дисплеем делала библиотека. Здесь обойдёмся без неё, чтобы увидеть, что
происходит на самом деле. Датчик --- это набор регистров (урок 22), а шина I\textsuperscript{2}C позволяет
записать число в регистр или прочитать несколько регистров подряд.

\begin{center}\small
\begin{tabular}{@{}lll@{}}
\toprule
Адрес регистра & Название & Назначение \\
\midrule
\code{0x75} & WHO\_AM\_I & «паспорт»: исправный датчик отвечает \code{0x68} \\
\code{0x6B} & PWR\_MGMT\_1 & питание: запись 0 выводит датчик из сна \\
\code{0x3B}--\code{0x40} & ACCEL\_X, Y, Z & ускорение: по два байта на ось, старший первым \\
\code{0x43}--\code{0x48} & GYRO\_X, Y, Z & скорость поворота: по два байта на ось \\
\bottomrule
\end{tabular}
\end{center}

Каждое измерение --- число со знаком из 16 бит, которое приходит двумя байтами. Собрать его обратно
помогает сдвиг: \code{(старший << 8) | младший}. В диапазоне по умолчанию 16384 единицы акселерометра
соответствуют $1g$, а 131 единица гироскопа --- одному градусу в секунду.

Модуль GY-521 подключается как любое устройство I\textsuperscript{2}C: VCC к 3V3, GND к земле, SDA к \pin{8},
SCL к \pin{9}. Его адрес --- \code{0x68}; проверьте это сканером из урока 7.

\begin{figure}[H]
\centering
\begin{circuitikz}[scale=0.95,transform shape]
  \node[blk,fill=blkcpu,minimum width=2.8cm,minimum height=3.8cm] (m) at (0,0) {};
  \node[blk,fill=blkper,minimum width=2.8cm,minimum height=3.8cm] (d) at (8,0) {};
  \node[font=\sffamily\small\bfseries] at (0,2.3) {ESP32-S3};
  \node[font=\sffamily\small\bfseries,align=center] at (8,2.45) {GY-521\\[-2pt]\scriptsize MPU6050};
  \foreach \y/\n in {1.5/3V3,0.5/{\pin{8}},-0.5/{\pin{9}},-1.5/GND}
    \node[font=\sffamily\scriptsize,anchor=east] at (1.4,\y) {\n};
  \foreach \y/\n in {1.5/VCC,0.5/SDA,-0.5/SCL,-1.5/GND}
    \node[font=\sffamily\scriptsize,anchor=west] at (6.6,\y) {\n};
  \draw[thick,esp]      (1.4,1.5)  -- (6.6,1.5);
  \draw[thick,accent]   (1.4,0.5)  -- (6.6,0.5);
  \draw[thick,okgreen]  (1.4,-0.5) -- (6.6,-0.5);
  \draw[thick,black!70] (1.4,-1.5) -- (6.6,-1.5);
  \node[font=\sffamily\scriptsize,text=accent,fill=white,inner sep=1pt] at (4,0.5) {данные};
  \node[font=\sffamily\scriptsize,text=okgreen,fill=white,inner sep=1pt] at (4,-0.5) {такт};
  \node[font=\sffamily\scriptsize,align=center,gray!40!black] at (8,-2.5) {выводы AD0, INT, XDA, XCL\\[-2pt] остаются свободными};
\end{circuitikz}
\caption{Подключение GY-521 по I\textsuperscript{2}C. Подтяжки шины и стабилизатор уже есть на модуле,
поэтому внешних деталей не нужно. При свободном AD0 адрес равен \code{0x68}}
\end{figure}

\subsection{Комплементарный фильтр}

Идея фильтра умещается в одну строку. Новый угол на 98 \% берётся из гироскопа (прошлый угол плюс
поворот за прошедшее время) и на 2 \% --- из акселерометра:
\[
  \varphi = 0{,}98\cdot(\varphi + \omega\,\Delta t) + 0{,}02\cdot\varphi_\text{акс} .
\]
Гироскоп отвечает за быстрые изменения, а маленькая постоянная добавка от акселерометра не даёт углу
уплыть.

\begin{figure}[H]
\centering
\begin{tikzpicture}
\begin{axis}[
  width=0.78\textwidth, height=5.4cm,
  xlabel={время, с}, ylabel={угол, °},
  xmin=0, xmax=6, ymin=-8, ymax=42,
  grid=major, grid style={gray!25},
  legend pos=south east, legend style={font=\scriptsize}, legend cell align=left,
  tick label style={font=\scriptsize}, label style={font=\small},
]
  \addplot[thin,accent,domain=0:6,samples=220] {30/(1+exp(-6*(x-1.5))) + 3*sin(deg(31*x)) + 2*sin(deg(53*x+1))};
  \addlegendentry{только акселерометр: шум}
  \addplot[thick,warn,dashed,domain=0:6,samples=60] {30/(1+exp(-6*(x-1.5))) + 1.6*x};
  \addlegendentry{только гироскоп: дрейф}
  \addplot[very thick,esp,domain=0:6,samples=80] {30/(1+exp(-5*(x-1.55)))};
  \addlegendentry{комплементарный фильтр}
\end{axis}
\end{tikzpicture}
\caption{Плату наклонили на 30° (иллюстрация). Акселерометр шумит, угол по гироскопу уплывает, фильтр
берёт лучшее от обоих}
\end{figure}

Программа растёт в три шага.

\subsubsection{Шаг 1. Связь с датчиком}

Сначала убедимся, что датчик отвечает. Пишем две функции: \code{writeReg()} записывает число в регистр,
\code{readRegs()} читает несколько регистров подряд. В \code{setup()} читаем «паспорт» и будим датчик.

\codecap{Шаг 1. WHO\_AM\_I и пробуждение}
\codeinput{l26_mpu6050_s1}

\begin{itemize}[leftmargin=*]
  \item \code{Wire.endTransmission(false)} --- «я ещё не закончил»: шина остаётся занятой до чтения;
  \item исправный датчик отвечает \code{WHO\_AM\_I = 0x68}; если ответа нет, проверьте провода;
  \item запись нуля в \code{PWR\_MGMT\_1} выводит датчик из сна: пока он спит, все значения нулевые.
\end{itemize}

\subsubsection{Шаг 2. Угол по акселерометру}

Теперь читаем 14 байт подряд, начиная с регистра \code{0x3B}: ускорения, температуру и скорости поворота.
Берём ускорения по Y и Z, собираем из пар байтов числа со знаком и вычисляем угол по формуле с рисунка.

\codecap{Шаг 2. Наклон только по акселерометру}
\codeinput{l26_mpu6050_s2}

\begin{itemize}[leftmargin=*]
  \item \code{(b[2] << 8) | b[3]} --- старший байт сдвигается влево и склеивается с младшим;
  \item \code{(int16\_t)} --- число становится знаковым, поэтому отрицательные значения не теряются;
  \item \code{atan2(ay, az)} --- арктангенс отношения, результат в радианах; умножение на 57,2958 даёт градусы.
\end{itemize}

В Serial Plotter линия \code{acc} будет вздрагивать от любой вибрации --- именно поэтому нужен второй датчик.

\subsubsection{Шаг 3. Гироскоп и комплементарный фильтр}

В последнем шаге добавляем скорость поворота \code{gx} и время между измерениями \code{dt}. Фильтр из формулы
выше умещается в одну строку.

\codecap{Шаг 3. Чтение MPU6050 через регистры и комплементарный фильтр}
\codeinput{l26_mpu6050}

\begin{itemize}[leftmargin=*]
  \item \code{dt} --- сколько секунд прошло с прошлого замера; без него нельзя перевести скорость в угол;
  \item \code{gx / 131.0} --- скорость поворота в градусах в секунду (131 единица на градус);
  \item \code{0.98 * (\ldots) + 0.02 * accAngle} --- сам фильтр.
\end{itemize}

Откройте \textsf{Tools → Serial Plotter}, наклоняйте плату и постучите по столу: линия \code{acc} будет
вздрагивать, а \code{angle} останется спокойной.

\begin{tip}
Коэффициенты 0,98 и 0,02 --- не закон. Чем больше доля гироскопа, тем спокойнее, но медленнее
исправляется дрейф. Попробуйте 0,90 и 0,995 и сравните графики.
\end{tip}

\subsection{Самостоятельный опыт: электронный уровень}

Соедините датчик с OLED-дисплеем (оба устройства на одной шине I\textsuperscript{2}C). Покажите угол
крупными цифрами и нарисуйте «пузырёк», который смещается при наклоне. Если угол меньше 1°, зажигайте
зелёный светодиод.

\subsection{То же на ESPUI: график наклона}

График комплементарного фильтра в этом уроке был иллюстрацией; теперь
можно получить настоящий. Плата считает угол комплементарным фильтром, а телефон показывает
число по акселерометру, число после фильтра и ленту последних значений. Покачайте плату вокруг оси X:
график плавный, а «сырое» число заметно дрожит.

Датчик подключён так же, как раньше: SDA к GPIO8, SCL к GPIO9. Сеть называется \code{S3-Tilt}
(пароль \code{lesson-39-demo}), адрес панели --- \code{http://192.168.4.1/}. Строим панель в три шага.

\subsubsection{Шаг 1. График на пробу}

Сначала проверим саму панель, не трогая датчик: подпись и график, в которые \code{loop()} четыре раза
в секунду отправляет «синусоиду». Если график побежал --- панель в порядке.

\codecap{Шаг 1. Подпись и график с пробными данными}
\codeinput[\footnotesize]{l26_espui_tilt_s1}

\begin{itemize}[leftmargin=*]
  \item \code{ESPUI.graph(\ldots)} --- элемент-график; точки добавляет \code{addGraphPoint()};
  \item показ раз в 250 мс: телефону не нужно больше, а сеть не перегружается.
\end{itemize}

\subsubsection{Шаг 2. Настоящий угол по акселерометру}

Заменяем «синусоиду» углом из датчика. Функции \code{writeReg()} и \code{readRegs()} те же, что в основной
программе урока. Если датчик не отвечает, на странице появляется «НЕ ОТВЕЧАЕТ», и программа не выдумывает
нулевые углы.

\codecap{Шаг 2. Угол по акселерометру на графике}
\codeinput[\footnotesize]{l26_espui_tilt_s2}

\begin{itemize}[leftmargin=*]
  \item подпись «Датчик» сообщает, ответил ли датчик на \code{WHO\_AM\_I};
  \item \code{atan2(ay, az)} --- тот же расчёт угла, что в шаге 2 основной программы;
  \item покачайте плату: график и число дрожат при каждом толчке.
\end{itemize}

\subsubsection{Шаг 3. Фильтр}

Добавляем гироскоп и комплементарный фильтр. Теперь на странице два числа: «сырое» по акселерометру и
после фильтра, а на график идёт отфильтрованное.

\codecap{Шаг 3. Угол по акселерометру и после фильтра}
\codeinput[\footnotesize]{l26_espui_tilt}

\begin{trouble}
\begin{itemize}[leftmargin=*]
  \item \textbf{Датчик не отвечает} --- проверьте адрес сканером: при выводе AD0, подключённом к питанию, адрес равен \code{0x69}.
  \item \textbf{Все значения равны нулю} --- датчик спит: не записан 0 в регистр \code{PWR\_MGMT\_1}.
  \item \textbf{Угол на ровном столе не равен нулю} --- у каждого экземпляра свой сдвиг; измерьте его при
        запуске и вычитайте (калибровка нуля).
\end{itemize}
\end{trouble}

\begin{summary}
\begin{itemize}[leftmargin=*]
  \item Любой датчик I\textsuperscript{2}C --- это регистры: в одни пишут настройки, из других читают результат.
  \item Акселерометр даёт наклон по силе тяжести, гироскоп --- скорость поворота.
  \item Комплементарный фильтр смешивает оба показания и убирает и шум, и дрейф.
\end{itemize}
\end{summary}

\begin{selfcheck}
\begin{enumerate}
  \item Как из двух прочитанных байтов получить 16-битное число со знаком?
  \item Почему по одному гироскопу нельзя долго измерять угол?
  \item Что изменится, если в фильтре заменить коэффициенты 0,98 и 0,02 на 0,5 и 0,5?
\end{enumerate}
\end{selfcheck}

\begin{tasks}
\begin{enumerate}
\item \lv{1} Выведите в Serial Plotter ускорения по трём осям и найдите, какая ось куда направлена.
\item \lv{1} Положите плату на стол и проверьте, что суммарное ускорение близко к $1g$.
\item \lv{2} Добавьте калибровку: первые две секунды плата неподвижна, программа измеряет сдвиг гироскопа и вычитает его.
\item \lv{2} Рассчитайте угол наклона только по акселерометру и выведите его в градусах.
\item \lv{3} Сделайте детектор встряхивания: если полное ускорение превышает $2g$, плата подаёт звуковой сигнал.
\item \lv{4} Сравните угол по акселерометру, по гироскопу и по комплементарному фильтру на графике при быстром покачивании платы.
\end{enumerate}
\end{tasks}
